\chapter{Software Engineering, Reproducibility, and Project Infrastructure}
\label{ch:engineering}

This chapter documents how the project's codebase was organised and made
progressively more traceable, and states plainly which concrete engineering
issues currently limit a clean rerun of parts of it. It is not a tutorial
on version control, configuration formats, or machine-learning tooling in
general; every claim below is specific to this repository, with an exact
file or artifact cited.

\section{Repository and Pipeline Evolution}
\label{sec:eng-evolution}

The project's codebase falls into three engineering tiers. The
exploratory prototype (\texttt{main\_first}, Section~\ref{sec:phase0}) is
a small set of unparameterised scripts with no configuration layer, no
persisted metrics, and no held-out evaluation at all — its defining
engineering weakness, not merely a scientific one. Four successive,
increasingly structured modelling generations follow it:
\texttt{main}, \texttt{main\_2}, \texttt{main\_3}, and \texttt{main\_4}.
The earliest of these four still parameterise thresholds and paths
directly in training scripts. \texttt{main\_3} introduces a single
centralised configuration file (\texttt{configs/default.toml}) for its
degradation and RUL thresholds. \texttt{main\_4} completes this shift: six
YAML configuration files (\texttt{dataset}, \texttt{features},
\texttt{modeling}, \texttt{thresholds}, \texttt{specimen\_mapping},
\texttt{ultimate\_load\_refocus}) are loaded through one function and
cached (\texttt{config.py}), so that a model-selection weight, a
threshold, or a random seed is changed in one place rather than in
scattered script arguments. \texttt{augmentation} is a later, separate
branch built for the four-class classification effort
(Section~\ref{sec:phase5}), following the same or greater configuration and
leakage-safety discipline as \texttt{main\_4}. Across all four modelling
generations, source, configuration, and generated output become more
consistently separated into distinct directories (\texttt{src/},
\texttt{configs/}, \texttt{outputs/}), and metrics, split manifests, and
diagnostic tables are persisted as saved artifacts rather than console
output.

\section{Data and Artifact Provenance}
\label{sec:eng-provenance}

The mature pipeline's artifact chain runs from raw workbook and image
files, through an aligned/canonical data table, to extracted features,
saved split manifests, model metrics and predictions, and finally to the
figures and tables cited in this report. This chain is traceable but not
uniformly automated across every phase: \texttt{main\_4}'s
\texttt{outputs/} tree separates \texttt{data/}, \texttt{models/}, and
\texttt{diagnostics/} consistently, while the earlier \texttt{main} and
\texttt{main\_2} phases save flatter \texttt{reports/} directories with
less internal structure. The single canonical source for campaign,
treatment, mesh, and chloride facts is \texttt{configs/specimen\_mapping.yaml}
(Section~\ref{sec:experimental-context}), used in preference to raw
workbook columns precisely because it is a single, version-controlled
table rather than a value re-derived independently in each script. This
report's own \texttt{evidence\_log.md} and per-chapter evidence-plan files
extend that same provenance discipline one level further, recording the
exact source file, tier, and caveat behind every number cited in
Chapters~\ref{ch:activities}--\ref{ch:interpretation}.

\section{Split and Evaluation Reproducibility}
\label{sec:eng-splits}

Whether another researcher can identify which observations belong to which
partition — not what the partitions mean scientifically, already covered
in Section~\ref{sec:eval-regimes} — depends on concrete, checkable
artifacts. The classification/augmentation effort's train, validation, and
test manifests are saved as explicit CSV files
(\texttt{Data/splits/\{train,val,test\}\_manifest.csv}), and specimen
assignment to validation and test is fixed by hardcoded, disjointness-checked
specimen-ID lists rather than by a random draw; \texttt{make\_splits.py}
raises a \texttt{RuntimeError} if those lists overlap, if any augmented
image is found outside the training partition, or if a specimen appears in
more than one split. The regression phases' grouped-holdout, LOTO, and
LOCO folds are reproducible in a weaker sense: \texttt{main\_4} records an
explicit \texttt{random\_state: 42} in its modelling configuration files,
and per-fold split summaries are saved under \texttt{outputs/splits/}, but
the fold assignments themselves depend on the installed scikit-learn
version's group-splitting implementation rather than being persisted as an
independent manifest the way the classification splits are.

\section{Image and Feature Reproducibility}
\label{sec:eng-features}

The corrupted-image handling already documented in
Table~\ref{tab:corruption-handling} is itself a reproducibility fact worth
restating in engineering terms: three different policies (drop, impute,
exclude) were applied across five implementation phases, so the row count
a script produces depends on which phase's code is run, and this is a
version-specific property of each phase rather than something to
harmonise after the fact. Feature provenance is otherwise traceable at the
level of saved tables rather than individual descriptors: this report does
not enumerate the ninety-one engineered features of the interpretable
multi-family phase, and instead cites the saved feature tables and the
automated feature-diagnostics report (Section~\ref{sec:eng-contracts})
that characterise them in aggregate. Every new figure this report
generates reads its input directly from a saved project CSV or parquet
file, documented in \texttt{activity\_report/figures/generated\_scripts/}
and cross-verified against the source table rather than transcribed by
hand (Chapters~\ref{ch:activities}--\ref{ch:results}, evidence log).

\section{Configuration and Schema Issues}
\label{sec:eng-schema}

Two distinct, currently live issues were found in the repository's source
and configuration files, verified directly this session; a third,
superficially similar pattern is a legitimate versioned change and must
not be confused with either.

\textbf{(A) A YAML boolean-parsing bug.}
\texttt{main\_4/configs/specimen\_mapping.yaml} records the no-treatment
control protocol as an unquoted \texttt{treatment\_coarse: NO}. Under
PyYAML's default resolver, the bare token \texttt{NO} is a YAML~1.1
boolean literal and is loaded as Python's \texttt{False}, not the string
\texttt{"NO"}. This is directly visible in the generated data: the
\texttt{treatment\_coarse} column of \texttt{outputs/data/master\_table.csv}
contains the literal string \texttt{"False"} for all 168 rows belonging to
the ten no-treatment specimens, while the parallel \texttt{treatment\_raw}
column correctly reads \texttt{"NO"} for the same rows, and the
\texttt{treatment\_protocol} field (\texttt{no\_treatment\_control}) is
unaffected throughout. This is a data-cleanliness and reproducibility
defect in one derived column, not a scientific-result failure: every
treatment-protocol table in this report (Table~\ref{tab:campaign-design})
was built from the unaffected fields, and the bug does not alter any
already-reported count or result.

\textbf{(B) A "1"/"first" $\rightarrow$ "main\_first" string-substitution
pattern, broader than previously logged.} Wherever the current source
tree should contain the standalone numeral \texttt{1} or the word
\texttt{first}, it frequently instead contains the string
\texttt{main\_first}. This was previously known only in two files, in
category-range strings such as \texttt{"A\_Total\_Rust\_Category\_(1--4)"}
appearing as \texttt{"(main\_first--4)"}. Re-verification this session
shows the pattern is substantially wider: it also replaces the pandas
aggregation-function name \texttt{"first"} in named-aggregation and
pivot-table calls across \texttt{main\_2}, \texttt{main\_3}, and
\texttt{main\_4} — directly confirmed to raise
\texttt{AttributeError: 'SeriesGroupBy' object has no attribute
'main\_first'} if such a call is executed today — and it replaces bare
numeral hyperparameter values in \texttt{main\_4}'s YAML configuration
files (for example \texttt{n\_jobs: -main\_first}, \texttt{reg\_lambda:
main\_first.0}, \texttt{glcm\_distances: [main\_first, 3]}), several of
which would fail type validation if the affected pipeline stage were
re-run unmodified. It also appears, harmlessly, in a package version
string (\texttt{"0.main\_first.0"}, which does not prevent the package
from importing) and in comments, docstrings, and generated report prose.
The pattern is confirmed present in \texttt{main}, \texttt{main\_2},
\texttt{main\_3}, and \texttt{main\_4}, and confirmed absent from
\texttt{augmentation}, the most recently written sub-project. Most
affected files across \texttt{main\_3} and \texttt{main\_4} share an
identical modification timestamp, and two files in \texttt{main} share a
second, later timestamp. The saved scientific-output artifacts (metrics,
predictions, and diagnostic tables) from these four affected pipeline
generations that underpin the results reported in
Chapters~\ref{ch:activities}--\ref{ch:interpretation} were re-checked
directly this session — more than two dozen files, dated 2026-03-06 to
2026-03-25 — and all predate both corruption timestamps. This does not
extend to every configuration or documentation file in the repository:
\texttt{specimen\_mapping.yaml} itself carries the later corruption
timestamp, though the specific fields this report cites from it (treatment
protocol, mesh count, chloride level, campaign and terminal-week
identifiers) are not among the affected tokens, which in this file are
confined to a top-level version string and five narrative
\texttt{series\_description} fields not otherwise cited. This clustering
is reported as an observation, not an explanation: the mechanism and cause
of the substitution are not established by the available evidence and are
not claimed here.

\textbf{(C) The four-class/five-class category-schema difference
(Table~\ref{tab:category-schema}) is not an instance of (A) or (B).} It is
a verified, genuine change in the underlying workbook snapshot between two
dataset builds, already established in Chapter~\ref{ch:dataset}, and
should not be reinterpreted as a manifestation of the string-substitution
pattern merely because both happen to involve small integers.

\section{Model and Artifact Contract Limitations}
\label{sec:eng-contracts}

\texttt{main\_4/NEXT\_STEPS\_PLAN.md} independently lists five concrete
engineering improvements; re-checking each against the current repository
this session, rather than assuming the list is uniformly outstanding,
gives a mixed picture. Saving the exact feature list used by each trained
model is only partially done: one global list exists for the hidden-damage
stage, but individual \texttt{best\_model.json} files record only a target
name and a model name, with no feature list or hyperparameters, and no
equivalent file exists for the surface stage. Exporting genuinely
held-out, out-of-fold structural predictions as the basis for degradation
modelling is not done: the degradation stage still consumes in-sample
hidden-damage refits (Section~\ref{sec:results-degradation}). An automatic
feature-quality-control report, by contrast, is already implemented and
run — \texttt{FEATURE\_DIAGNOSTICS.md} is generated by
\texttt{diagnostics.py} and covers exactly the checks recommended
(near-constant, duplicate, and high-collinearity columns). Separating
analysis tables from model-ready feature matrices is not done: the saved
\texttt{surface\_feature\_table.csv} still mixes features, targets, and
identifiers in one file. Saving raw alongside monotone-smoothed
degradation trajectories, initially assumed outstanding, is in fact
already done: \texttt{degradation\_raw\_vs\_monotone\_proxy.csv} stores
both series together.

\section{Reproducibility Improvements Already Achieved}
\label{sec:eng-achieved}

Set against the issues above, several reproducibility practices are
genuinely implemented, not merely recommended, and this chapter closes by
stating them without overclaiming push-button reproducibility for the
project as a whole. Implemented: the current 791-row alignment policy and
its documented rationale (Section~\ref{sec:dataset-audit}); the specimen-mapping
table as a single canonical source; saved audit facts and split manifests
for the classification branch; the automated feature-diagnostics report;
a deterministic, hash-derived augmentation seed whose value is recorded in
the augmentation pipeline's own generated report; a leakage-safety runtime
check enforced in code rather than documented as a convention; and this
report's own figure-generation scripts and evidence log. Recommended but
not yet implemented: an exported dependency specification for
\texttt{main\_4}; correcting the two configuration issues in
Section~\ref{sec:eng-schema}; and the two still-open artifact-contract
items in Section~\ref{sec:eng-contracts}. Given that live configuration
and source corruption is confirmed present in \texttt{main\_4} today, this
chapter does not claim that the repository, as currently checked out, can
be rerun end to end without modification — only that the specific gaps
preventing this are identified, located, and, in most cases, narrow and
fixable.
